% Evans 12.1, printed pp.617--620; English PDF pp.634--637.
% Chinese section start: printed p.437 / scan p.455.
\section{Introduction}
\label{sec:12-1}
\subsection{The equation and its conserved energy}
Write $\PDEWave=\partial_t^2-\Delta_x$. The main equation is
\begin{equation}\label{eq:12-semi}
 \PDEWave u+f(u)=0,\qquad u(0)=g,\quad u_t(0)=h
 \quad\hbox{on }\R^n.
\end{equation}
We also study $\PDEWave u+N(Du,u_t,u)=0$, with $N(0,0,0)=0$.
Its highest-order part is still the constant wave operator.
Although Evans calls it mildly quasilinear, its nonlinearity involves
only lower derivatives; equations with solution-dependent principal
coefficients require additional analysis. Throughout $D=D_x$.
For a scalar $f$ set $F(z)=\int_0^z f(a)\,da$, so $F(0)=0$.
\PDEAnnotation{非线性项的符号}{本章把 $f$ 放在左侧。
$f(u)=|u|^{p-1}u$ 给非负势能，$f(u)=-|u|^{p-1}u$ 给负势能。
后者才可能满足负能量爆破条件；不能只比较幂次而忽略符号。}

\begin{theorem}[Energy; source \S12.1.1, Theorem 1]
For a smooth solution with support in one common compact spatial
set on each finite time strip, or sufficient decay to justify the
same integrations,
\begin{equation}\label{eq:12-energy}
 E(t)=\int_{\R^n}\left\{\frac12(u_t^2+|Du|^2)+F(u)\right\}dx=E(0).
\end{equation}
\end{theorem}
\begin{expandedproof}[local identity before integration]
Put $e=\frac12(u_t^2+|Du|^2)+F(u)$. The product rule gives
\[
 e_t-\operatorname{div}(u_tDu)=u_t(\PDEWave u+f(u))=0.
\]
Integrate on a ball containing the support over a fixed time strip;
its boundary term is zero. Under decay hypotheses use cutoffs with
radius tending to infinity and check that $u_tDu$ is integrable.
This also explains why control of both kinetic and gradient energy
is sufficient to eliminate the flux at infinity. For weak solutions
the product $u_t\PDEWave u$ may be undefined; conservation must then be
established by approximation, not by this formal multiplication.
\end{expandedproof}

\subsection{Backward cones and the exact flux}
For apex $(x_0,T)$ put
\[
 \begin{aligned}
 B_t&=B(x_0,T-t),\\
 K_{s,\tau}&=\{(x,t):s<t<\tau,\ x\in B_t\},\\
 \Gamma_{s,\tau}&=\{(x,t):s<t<\tau,\ |x-x_0|=T-t\}.
 \end{aligned}
\]
On the mantle let $\nu=(x-x_0)/|x-x_0|$. Write
$E_K(t)=\int_{B_t}e(x,t)\,dx$.
\PDEAnnotation{锥面上的两种测度}{锥面 $t=T-|x-x_0|$ 的诱导测度
为 $\sqrt2\,dx$，分时切片则为 $\sqrt2\,dS_xdt$。
因此式中的 $1/\sqrt2$ 是几何因子。
临界点可能奇异时，顶端剩余能量不能预先设为零。}
\begin{theorem}[Cone flux; source \S12.1.2, Theorem 2]
For a solution smooth up to every truncated cone,
\begin{equation}\label{eq:12-cone-flux}
 E_K(s)-E_K(\tau)=\frac1{\sqrt2}\int_{\Gamma_{s,\tau}}
 \left\{\tfrac12|u_t\nu-Du|^2+F(u)\right\}dS.
\end{equation}
If $F\ge0$, the local cone energy is nonincreasing. If its initial
section has $g=h=0$, then $u=0$ in that cone.
\end{theorem}
\begin{expandedproof}[moving boundary and the surface factor]
Reynolds' formula for the ball shrinking at unit speed gives
\[
 E_K'(t)=\int_{\partial B_t}(u_t\partial_\nu u-e)\,dS_x
 =-\int_{\partial B_t}\{\tfrac12|u_t\nu-Du|^2+F(u)\}\,dS_x.
\]
The square is $u_t^2+|Du|^2-2u_t\partial_\nu u$; its tangential
part must not be dropped. Integrating in time proves the result,
because the space-time mantle measure satisfies
$dS=\sqrt2\,dS_xdt$. If the initial local energy is zero, positivity
forces $u_t=Du=0$ in every smaller section. Integrating along vertical
segments from the zero initial data gives $u=0$ there, even if $F$
has other zeroes. When $u$ is smooth at the apex, $E_K(\tau)\to0$
as $\tau\uparrow T$. We retain this term whenever the apex might
be singular: smoothness only for $t<T$ does not justify dropping it.
\end{expandedproof}
\PDEBodyDiagram{12.1}{The shrinking sections and the null flux}
{A spatial section of the backward unit-speed cone has radius $T-t$.
Energy leaves through the mantle with density
$\frac12|u_t\nu-Du|^2+F(u)$, measured using $dS/\sqrt2$.
The top section at $\tau<T$ is retained. The drawing is a one-space-dimensional
section of the $n$-dimensional cone, not a reduction of the PDE.}
{ch12-cone-flux}

\subsection{Finite propagation without a positive potential}
\begin{theorem}[Domain of dependence; source \S12.1.2, Theorem 3]
Let $N$ be locally Lipschitz, $N(0)=0$, and let two classical solutions
of $\PDEWave u+N(Du,u_t,u)=0$ have bounded first-order state on a compact
backward cone. If their two initial data agree on its base, they
agree inside the cone. In particular compactly supported data yield
support in $B(0,R+t)$ while a classical solution exists.
\end{theorem}
\begin{expandedproof}[difference energy and local Lipschitz constants]
For $w=u-v$, write $Z_u=(Du,u_t,u)$ and define $Z_v$ similarly.
The local Lipschitz bound is
\[
 |N(Z_u)-N(Z_v)|\le L(|Dw|+|w_t|+|w|),
\]
where $L$ depends only on a bound for the two states in this cone.
Add $w^2/2$ to the usual wave energy and integrate on its shrinking
sections. The lateral boundary contribution is
$-\frac12|w_t\nu-Dw|^2-w^2/2\le0$.
The interior contribution is $w_t\{w-(N(u)-N(v))\}$ and is at most
$C_L(w_t^2+|Dw|^2+w^2)$. Thus $Y'(t)\le C_LY(t)$ and $Y(0)=0$;
Gronwall gives $Y=0$. Comparison with zero uses $N(0)=0$.
For energy solutions with globally Lipschitz $N$, the same inequality
is obtained by time/spatial regularization and smooth cone cutoffs,
then passage to the limit in $L^2$. Thus it is also available for the
local-data gluing in the next section. A cone is characteristic for
the fixed principal wave operator; adding lower-order derivative
terms does not change its speed.
\end{expandedproof}
